\documentclass[conference]{IEEEtran}
\IEEEoverridecommandlockouts
\usepackage{cite}
\usepackage{amsmath,amssymb,amsfonts}
\usepackage{algorithmic}
\usepackage{graphicx}
\usepackage{textcomp}
\usepackage{xcolor}
\usepackage{booktabs}
\usepackage{url}
\def\BibTeX{{\rm B\kern-.05em{\sc i\kern-.025em b}\kern-.08em
    T\kern-.1667em\lower.7ex\hbox{E}\kern-.125emX}}

% ============================================================================
%  MMS 2025/26 — Project for Additional Points
%  Topic #7: OSCAR — One-Step Diffusion Codec Across Multiple Bit-rates
%  Authors: Eldar Buzadzic, Edi Djelmo
%
%  Section plan (8 sections, ~10 pages IEEE two-column):
%    I.    Introduction                       (problem statement)            ~1p
%    II.   The OSCAR Algorithm                (algorithm explanation, math)  ~2-2.5p
%    III.  Implementation on AMD/ROCm         (what we did)                  ~1.5-2p
%    IV.   Evaluation Methodology             (metrics + dataset)            ~1p
%    V.    Problems Faced and Resolutions     (failure modes + fixes)        ~1.5-2p
%    VI.   Results and Discussion             (numbers + interpretation)     ~1.5-2p
%    VII.  Possible Improvements              (future work)                  ~0.5-0.75p
%    VIII. Conclusion                                                        ~0.25p
% ============================================================================

\begin{document}

\title{Reproducing OSCAR on ImageNet via AMD ROCm:\\ A One-Step Multi-Rate Diffusion Image Codec}

\author{\IEEEauthorblockN{Eldar Buzad\v{z}i\'{c}}
\IEEEauthorblockA{\textit{Faculty of Electrical Engineering} \\
\textit{University of Sarajevo}\\
Sarajevo, Bosnia and Herzegovina \\
ebuzadzic1@etf.unsa.ba}
\and
\IEEEauthorblockN{Edi \DJ elmo}
\IEEEauthorblockA{\textit{Faculty of Electrical Engineering} \\
\textit{University of Sarajevo}\\
Sarajevo, Bosnia and Herzegovina \\
edelmo1@etf.unsa.ba}
}

\maketitle

\begin{abstract}
OSCAR is a recently published image codec that uses a single denoising
step of a latent diffusion model to reconstruct images at multiple
bit-rates, mapping each rate to a pseudo-timestep on the diffusion
schedule. We reproduce OSCAR on a 5000-image seed-controlled subset of
the ImageNet-1k validation split, running on a consumer AMD Radeon RX
9070 XT (RDNA 4) via ROCm 7.2.1. The reproduction required porting the
CUDA-only upstream codebase to AMD's Windows PyTorch fork, sourcing
weights from a community mirror after Stability AI delisted the
canonical Stable Diffusion 2.1 repository, adapting the input pipeline
to ImageNet's varied aspect ratios, and resolving a small number of
numerical and operational pitfalls specific to the Windows ROCm
environment. The released checkpoint, evaluated under the original
paper's seven-metric protocol (PSNR, MS-SSIM, DISTS, LPIPS, MUSIQ,
CLIP-IQA, FID), reproduces the monotone rate-distortion behaviour
reported by the authors, with FID dropping from 110.64 at 0.0019\,bpp
to 5.78 at 0.125\,bpp across a four hour twenty seven minute
end-to-end run.
\end{abstract}

\begin{IEEEkeywords}
image compression, latent diffusion, OSCAR, ImageNet, ROCm, RDNA 4
\end{IEEEkeywords}

% ----------------------------------------------------------------------------
\section{Introduction}
\label{sec:intro}

Lossy image compression has long been guided by rate-distortion objectives that
optimize signal-fidelity metrics such as the peak signal-to-noise ratio (PSNR)
or the structural similarity index (SSIM). At low bit-rates, however, these
objectives produce reconstructions that are close in pixel space yet exhibit
visible blocking, blurring, and ringing artifacts that a human observer
immediately recognizes as degraded. \emph{Generative} codecs reframe the
problem: rather than minimizing pixel-wise error, they sample from a learned
distribution conditioned on a compact code, trading some signal fidelity for
a dramatic improvement in perceptual realism~\cite{hific}.

Latent diffusion models~\cite{ldm} have emerged as a particularly strong prior
for generative compression. The diffusion forward process is, by construction,
a model of how information is lost as latents are progressively noised; the
reverse process is a powerful learned reconstructor of clean images from noisy
ones. Diffusion-based codecs exploit this by encoding an image into a quantized
latent code that the diffusion model treats as a noisy intermediate state and
denoises back into a high-fidelity reconstruction~\cite{perco,codiff}.

Two practical limitations have, until recently, kept diffusion-based codecs
costly. First, the reverse process typically requires many forward passes
($T \in [20, 1000]$) through a large U-Net, putting inference one to two
orders of magnitude behind learned-CNN codecs. Second, each target bit-rate has
historically required a separately trained model, multiplying both training time
and storage. OSCAR~\cite{oscar} addresses both. Its key insight is that
compressed latents at different bit-rates correspond to noisy variants of the
original latent at different levels of distortion, so a single diffusion U-Net
can handle multiple rates by conditioning on a \emph{pseudo-timestep} mapped
from the bit-rate. Combined with a single denoising step at inference, OSCAR
becomes a one-pass, multi-rate codec.

In this work we reproduce OSCAR on the ImageNet-1k validation set~\cite{imagenet},
running entirely on consumer AMD hardware (Radeon RX 9070 XT, RDNA 4) via
ROCm 7.2.1. We make three contributions. First, we port the original
CUDA-only codebase to AMD's PyTorch fork on Windows, addressing several
silent-failure modes specific to that environment. Second, we evaluate the
released OSCAR checkpoint on a seed-controlled 5000-image subset of the
validation split, an order of magnitude larger than the largest set the
original paper reports on (CLIC2020, 428 images), and on a domain (natural
photographs at varied aspect ratios) that the original paper does not test.
Third, we report the full seven-metric battery used by the authors (PSNR,
MS-SSIM, DISTS, LPIPS, MUSIQ, CLIP-IQA, FID) across all eight of OSCAR's
bit-rate operating points. The resulting reconstructions show monotonic
rate-distortion behavior on every metric, with FID dropping from 110.6 at
$0.0019$\,bpp to $5.8$ at $0.125$\,bpp.

The remainder of this paper is organized as follows.
Section~\ref{sec:algorithm} details the OSCAR algorithm.
Section~\ref{sec:implementation} describes our implementation on AMD ROCm.
Section~\ref{sec:methodology} defines the evaluation methodology.
Section~\ref{sec:problems} documents the practical issues encountered and their
resolutions. Section~\ref{sec:results} presents the results.
Section~\ref{sec:improvements} discusses possible improvements, and
Section~\ref{sec:conclusion} concludes.

% ----------------------------------------------------------------------------
\section{The OSCAR Algorithm}
\label{sec:algorithm}

OSCAR is built on the latent diffusion paradigm and extends it with three
design choices that turn a multi-step generative model into a multi-rate,
one-step codec: a fixed mapping from bit-rate to a discrete \emph{pseudo-timestep},
a bank of per-rate hyper-encoders with vector quantization, and a magnitude
rescaling that lets a single diffusion U-Net treat compressed latents as drop-in
noisy intermediates. We unpack each of these in turn.

\subsection{Latent Diffusion Background}

Let $x \in \mathbb{R}^{3 \times H \times W}$ denote an input image. Following
Stable Diffusion 2.1~\cite{ldm}, OSCAR uses a frozen variational autoencoder
$(\mathcal{E}, \mathcal{D})$ that maps between pixel space and a 4-channel latent
space at one-eighth spatial resolution:
\begin{equation}
z_0 = s \cdot \mathcal{E}(x) \in \mathbb{R}^{4 \times H/8 \times W/8},
\label{eq:vae-encode}
\end{equation}
where $s = 0.18215$ is the VAE scaling factor. The diffusion forward process
applied to $z_0$ is a Markov chain that progressively adds Gaussian noise,
\begin{equation}
q(z_t \mid z_0) = \mathcal{N}\!\left(z_t;\, \sqrt{\bar{\alpha}_t}\, z_0,\, (1-\bar{\alpha}_t)\, \mathbf{I}\right),
\label{eq:forward}
\end{equation}
with $\bar{\alpha}_t = \prod_{i=1}^{t}\alpha_i$ the cumulative noise schedule of
the underlying variance-preserving DDPM~\cite{ddpm}. Reparameterizing,
\begin{equation}
z_t = \sqrt{\bar{\alpha}_t}\, z_0 + \sqrt{1-\bar{\alpha}_t}\, \epsilon, \quad \epsilon \sim \mathcal{N}(0, \mathbf{I}).
\label{eq:forward-repar}
\end{equation}
The SD-2.1 U-Net is trained under the v-prediction parameterization~\cite{vpred}.
Instead of predicting $\epsilon$ directly, the network outputs
$v = \sqrt{\bar{\alpha}_t}\,\epsilon - \sqrt{1-\bar{\alpha}_t}\, z_0$, from which
the clean latent is recoverable in closed form:
\begin{equation}
\hat{z}_0 = \sqrt{\bar{\alpha}_t}\, z_t - \sqrt{1-\bar{\alpha}_t}\, v_\theta(z_t, t, e).
\label{eq:v-to-x0}
\end{equation}
Equation~\eqref{eq:v-to-x0} is what OSCAR exploits to obtain a reconstruction
from a single forward pass.

\subsection{Mapping Bit-Rate to Pseudo-Timestep}

OSCAR's central insight is that a quantized latent at bit-rate $r$ can be
modeled \emph{as if} it were a noised version of the true latent at some
specific diffusion timestep $t(r)$, with higher distortion (lower bit-rate)
corresponding to a larger $t$. This collapses the family of ``one model per
rate'' into a single network conditioned on $t$. The authors define a fixed
monotone table of eight operating points (Table~\ref{tab:bpps}), with the
lowest rate of $0.0019$\,bpp mapped to $t = 500$ and the highest rate of
$0.125$\,bpp mapped to $t = 150$ on the standard $T = 1000$ schedule.

\begin{table}[h]
\caption{The eight OSCAR operating points: bit-rate, pseudo-timestep, and per-rate vector-quantizer configuration $(\text{cfg}_{ss}, \text{cfg}_{cs})$ controlling spatial size and codebook size.}
\label{tab:bpps}
\begin{center}
\begin{tabular}{rrrr}
\toprule
$r$ (bpp) & $t(r)$ & $\text{cfg}_{ss}$ & $\text{cfg}_{cs}$ \\
\midrule
0.0019 & 500 & 8  & 256 \\
0.0098 & 400 & 16 & 1024 \\
0.0313 & 310 & 32 & 256 \\
0.0430 & 280 & 32 & 2048 \\
0.0625 & 210 & 64 & 16 \\
0.0781 & 190 & 64 & 32 \\
0.0937 & 170 & 64 & 64 \\
0.1250 & 150 & 64 & 256 \\
\bottomrule
\end{tabular}
\end{center}
\end{table}

The table is hard-coded in the released implementation rather than learned. This
is a deliberate simplification: an eight-way table is enough to test the
central hypothesis and yields a clean rate-distortion curve at eight discrete
operating points. We discuss the option of a learned continuous mapping in
Section~\ref{sec:improvements}.

\subsection{Hyper-Encoder and Magnitude Rescaling}

For each bit-rate $r$, OSCAR trains a dedicated \emph{hyper-encoder}
$\mathcal{H}_r$ that compresses the VAE latent $z_0$ into a quantized code and
decodes it back into a latent-shaped tensor $\hat{z}$. The hyper-encoder is a
small convolutional network with three residual stages and a global attention
block, borrowed from PerCo~\cite{perco}; its forward pass can be summarized as
\begin{equation}
\hat{z} = \text{Proj}\!\left(\text{VQ}\!\left(\text{Backbone}(z_0)\right)\right),
\label{eq:hyper}
\end{equation}
where Backbone reduces $z_0$ to a feature map of spatial size $\text{cfg}_{ss}$,
the vector quantizer $\text{VQ}$ assigns each token to one of $\text{cfg}_{cs}$
codebook entries, and Proj upsamples the quantized tokens back to the original
latent shape. The pair $(\text{cfg}_{ss}, \text{cfg}_{cs})$ controls the rate: a
denser spatial grid and a larger codebook both increase the bit count per image.
The per-rate configurations are listed in Table~\ref{tab:bpps}.

A direct attempt to feed $\hat{z}$ into the U-Net at timestep $t(r)$ fails: the
vector quantizer changes the L2 magnitude of the latent in unpredictable ways,
while the U-Net expects inputs whose magnitude matches the noised latent
distribution at that timestep. OSCAR introduces a magnitude rescaling that
restores the original L2 norm,
\begin{equation}
\tilde{z} = \hat{z} \,\cdot\, \frac{\lVert z_0 \rVert_2}{\lVert \hat{z} \rVert_2},
\label{eq:rescale}
\end{equation}
with the norm computed along the spatial dimensions of each channel
independently. This single line of arithmetic is what allows the same fixed
U-Net to operate on the output of any of the eight hyper-encoders.

\subsection{Single-Step Reconstruction}

At inference, OSCAR replaces multi-step DDPM sampling with a single application
of the U-Net at the bit-rate's pseudo-timestep. For an image $x$ to be coded at
rate $r$ with $t = t(r)$:
\begin{align*}
z_0       &= s \cdot \mathcal{E}(x), \\
\tilde{z} &= \text{Rescale}\!\left(\mathcal{H}_r(z_0),\, z_0\right), \\
v         &= v_\theta(\tilde{z},\, t,\, e), \\
\hat{z}_0 &= \sqrt{\bar{\alpha}_t}\, \tilde{z} - \sqrt{1-\bar{\alpha}_t}\, v, \\
\hat{x}   &= \mathcal{D}(\hat{z}_0 / s),
\end{align*}
where $e \in \mathbb{R}^{77 \times 1024}$ is a \emph{learned} prompt embedding
that replaces the SD-2.1 text encoder. OSCAR stores $e$ as a single set of
parameters jointly with rank-16 LoRA~\cite{lora} adapters on the U-Net's
attention, projection, and convolutional layers; no text input is required at
coding time. The whole pipeline therefore costs one VAE encode, one small
hyper-encoder pass, one full U-Net pass, and one VAE decode, regardless of the
target bit-rate.

The model is fine-tuned in two stages. Stage~1 trains each hyper-encoder
$\mathcal{H}_r$ in isolation using a REPA-style cosine alignment between the
rescaled output and the original latent. Stage~2 jointly fine-tunes the LoRA
adapters on the U-Net, the eight hyper-encoders, and the learned embedding $e$
on a combined $L_1$ plus DISTS objective. We use the released stage-2
checkpoint as-is and do not retrain.

% ----------------------------------------------------------------------------
\section{Implementation on AMD ROCm}
\label{sec:implementation}

This section describes the reproduction implementation as it stands after the
porting and debugging effort. Section~\ref{sec:problems} documents the practical
issues that forced each design choice; here we report the final state.

\subsection{Hardware and Software Stack}

All experiments run on a single workstation built around an AMD Radeon
RX 9070 XT (16\,GB VRAM, RDNA 4 architecture, LLVM target \texttt{gfx1201}),
an AMD Ryzen 7 9800X3D 8-core CPU, and 32\,GB of system RAM, under Windows 11.
The choice of AMD hardware is deliberate: the original OSCAR codebase is
CUDA-only, and porting it to a consumer AMD card via ROCm is non-trivial in
ways we believe are worth documenting (Section~\ref{sec:problems}).

The Python environment uses Python 3.12 and consists of:
\begin{itemize}
  \item PyTorch \texttt{2.9.1+rocm7.2.1} (the \texttt{torch}, \texttt{torchvision},
        and \texttt{torchaudio} wheels), installed from AMD's official
        ROCm-on-Windows wheel index at \texttt{repo.radeon.com}. The
        accompanying ROCm SDK 7.2.1 ships as four \texttt{pip} packages
        (\texttt{rocm-sdk-core}, \texttt{rocm-sdk-devel},
        \texttt{rocm-sdk-libraries-custom}, \texttt{rocm}) totalling 1.4\,GB.
  \item \texttt{diffusers 0.25.0} and \texttt{transformers 4.37.2}, pinned to the
        versions OSCAR's vendored \texttt{autoencoder\_kl.py} and
        \texttt{unet\_2d\_condition.py} were derived from.
  \item \texttt{peft 0.9.0} (LoRA adapters on the U-Net), \texttt{compressai 1.2.6}
        (hyper-encoder convolution helpers), and
        \texttt{vector-quantize-pytorch 1.29.1} (the VQ codebook).
  \item \texttt{pyiqa 0.1.13} for the perceptual and no-reference quality
        metrics, plus the standard \texttt{Pillow}, \texttt{numpy},
        \texttt{huggingface\_hub}, and \texttt{datasets} stack.
\end{itemize}

Under PyTorch's ROCm port the AMD GPU is exposed as the \texttt{cuda} device,
so every \texttt{.to(``cuda'')} call in the upstream codebase works unchanged.
This is convenient but also misleading: not every CUDA-tested code path
exercises cleanly on ROCm, and Section~\ref{sec:problems} reports two such
cases. The \texttt{xformers} memory-efficient attention library, used by the
upstream training script, has no AMD build and is disabled; PyTorch's native
\texttt{scaled\_dot\_product\_attention} routes through ROCm's AOTriton backend
and is the fastest attention path available on this hardware.

\subsection{SD-2.1 Weight Acquisition}

The OSCAR checkpoint depends on the public Stable Diffusion 2.1 backbone for
its VAE and U-Net. Between the publication of OSCAR and this reproduction,
Stability AI removed the canonical \texttt{stabilityai/stable-diffusion-2-1}
repository from the HuggingFace Hub: \texttt{huggingface\_hub.snapshot\_download}
now returns HTTP 404 even with a valid authenticated token, and a directory
listing of the \texttt{stabilityai} organization confirms that none of the
SD-2.x family is currently hosted.

We reconstruct the weights from the community-hosted
\texttt{webui/stable-diffusion-2-1} repository, which retains the original
release as a 5.2\,GB FP32 \texttt{safetensors} file alongside its original
training-time YAML configuration. The diffusers single-file loader is
invoked with \texttt{prediction\_type=``v\_prediction''} and
\texttt{image\_size=768}, matching the canonical SD-2.1 release. Three string
replacements in \texttt{diffusers}'s \texttt{convert\_from\_ckpt.py}
redirect the hard-coded references to \texttt{stabilityai/stable-diffusion-2}
(the source of the CLIP-ViT-H text encoder configuration, also delisted) to a
still-available mirror at \texttt{sd-research/stable-diffusion-2-1-base}, whose
text-encoder configuration is identical to the original. The converted pipeline
is written to \texttt{model\_zoo/stable-diffusion-2-1/} in the standard
diffusers folder layout. The OSCAR fine-tuning checkpoint \texttt{oscar.pkl}
(394\,MB) is downloaded directly from its public repository at
\texttt{jinpeig/OSCAR}, which is not gated.

\subsection{Inference Pipeline}

The released \texttt{main\_test.py} is modified along four axes.

\textbf{Input geometry.}
The upstream pipeline trims each image to dimensions divisible by 128
(\texttt{w \% 128 = 0}, \texttt{h \% 128 = 0}) and feeds the result through the
model at native resolution. This is appropriate for the paper's evaluation sets
(Kodak, DIV2K-val, CLIC2020), all of which have consistent high-resolution
geometry. ImageNet validation images, by contrast, range from 60 to 3264 pixels
on the long side with diverse aspect ratios; under the upstream trim, small
images zero out and large images inflate runtime quadratically. We replace the
trim with an aspect-preserving resize to a short side of 384 pixels followed by
a center crop to $384 \times 384$,
\begin{equation*}
x' = \text{CenterCrop}_{384}\!\left(\text{Resize}_{\min=384}(x)\right),
\end{equation*}
which produces a fixed $48 \times 48$ latent for every input. This treatment
mirrors the paper's own protocol on CLIC2020 and DIV2K-val, both of which were
center-cropped to a fixed square.

\textbf{Mixed-precision execution.}
The upstream \texttt{main\_test.py} accepts a \texttt{mixed\_precision} CLI
option with allowed value \texttt{fp16} but never applies the resulting dtype
to the model; the option is, in effect, a no-op. We wrap the model forward
pass in a \texttt{torch.autocast} context,
\begin{verbatim}
with torch.no_grad(), torch.autocast(
        device_type="cuda",
        dtype=weight_dtype,
        enabled=(weight_dtype != torch.float32)):
    img_E, _, _ = model(lq, level)
img_E = img_E.float()
\end{verbatim}
which delivers a 1.7$\times$ throughput improvement on the 9070 XT while
leaving the model weights and OSCAR's internal hardcoded
\texttt{.to(torch.float32)} cast on the learned prompt embedding
(line 134 of \texttt{diffusion/oscar.py}) untouched. The trailing
\texttt{.float()} guarantees that downstream metric tensors are FP32.

\textbf{Per-metric robustness.}
Each of the seven \texttt{pyiqa} metric calls is wrapped in a helper that
returns \texttt{NaN} on exception:
\begin{verbatim}
def _safe_metric(name, fn, *args, **kwargs):
    try:
        v = fn(*args, **kwargs)
        return v.item() if hasattr(v, "item") else float(v)
    except Exception as e:
        print(f"[metric {name} failed: {e}]")
        return float("nan")
\end{verbatim}
and the per-bpp averages are computed by a NaN-skipping mean. A single image's
metric failure can no longer kill the run over all eight bit-rates
(Section~\ref{sec:problems}).

\textbf{DISTS on CPU.}
The DISTS metric, computed on CUDA in FP32, produces \texttt{NaN} for all
inputs on this ROCm build, but works correctly on CPU and on CUDA in FP64
(see Section~\ref{sec:problems} for the root cause). The metric is therefore
instantiated with \texttt{device=``cpu''} and its inputs are explicitly copied
to host memory at the call site: \texttt{dists\_metric(img\_E.cpu(),
img\_H.cpu())}. At roughly 80\,ms per call this adds about 55 minutes to a
5000-image $\times$ 8-bpp evaluation, an acceptable overhead.

\subsection{Unattended Execution}

A full evaluation run takes roughly 4.5 hours on the 9070 XT. To allow
overnight execution we (a) launch the eval as a fully detached Windows
process via PowerShell \texttt{Start-Process} with the
\texttt{WindowStyle Hidden} flag, with \texttt{stdout} and \texttt{stderr}
redirected to log files; (b) disable system sleep with
\texttt{powercfg /change standby-timeout-ac 0}, after a silent stall in an
early overnight attempt was traced to Windows entering \texttt{System Idle}
sleep mode 14 minutes into the run; and (c) initialize PyTorch with a fixed
seed of 42 and a stable image ordering (alphabetical, controlled by the
filename pattern of the sampling script described in
Section~\ref{sec:methodology}). Together these allow a single
\texttt{Start-Process} invocation to produce a deterministic
\texttt{results.csv} by morning.

% ----------------------------------------------------------------------------
\section{Evaluation Methodology}
\label{sec:methodology}

\subsection{Dataset}

We evaluate on a seed-controlled random subset of the ILSVRC2012 validation
split~\cite{imagenet}. The full validation set contains 50{,}000 labelled
images; we sample 5000 with \texttt{seed=42}, drawn from the ungated
HuggingFace mirror \texttt{Tsomaros/Imagenet-1k\_validation}, which hosts the
validation split as Apache Parquet shards. We decode each sampled record back
to a JPEG file and save the 5000 images to \texttt{data/imagenet\_val/} for
reproducibility. A manifest file records the source index of each sampled
image alongside the sampling seed and protocol.

The 10\,\% subset size is well above the paper's own evaluation budget: OSCAR
reports on Kodak (24 images), DIV2K-val (100 images), and CLIC2020 (428
images), so 5000 is more than an order of magnitude above the authors'
largest set. All images are processed at $384 \times 384$ after the input
pipeline described in Section~\ref{sec:implementation}, matching the paper's
fixed-square center-crop convention on its high-resolution test sets.

\subsection{Metrics}

We report the same seven metrics as the original paper, all computed through
the \texttt{pyiqa} library~\cite{pyiqa}. The mix spans full-reference and
no-reference, pixel-level and perceptual, and per-image and distributional
quality measures.

\textbf{PSNR.} Peak signal-to-noise ratio is the canonical pixel-level
distortion metric, expressed in dB. Higher is better. PSNR rewards exact
pixel matches and is known to penalize plausible-but-different generative
reconstructions; it is included for direct comparability with the wider
compression literature.

\textbf{MS-SSIM.} Multi-scale structural similarity~\cite{msssim} measures
local correlations of luminance, contrast, and structure at several
downsampled scales, returning a value in $[0, 1]$. Higher is better.
MS-SSIM tracks human-perceived quality better than PSNR at moderate to high
bit-rates.

\textbf{LPIPS.} Learned Perceptual Image Patch Similarity~\cite{lpips}
computes the $L_2$ distance between deep features of $x$ and $\hat{x}$
extracted from a pretrained VGG-16. Lower is better. LPIPS captures
perceptual differences that PSNR misses, such as small spatial misalignments
that look identical to a human observer.

\textbf{DISTS.} Deep Image Structure and Texture Similarity~\cite{dists}
extends LPIPS by separately modelling structural and textural similarity
through statistics over VGG features. Lower is better. DISTS is particularly
appropriate for generative codecs, whose reconstructions resynthesize
plausible textures rather than reproducing them pixel for pixel.

\textbf{MUSIQ.} The Multi-scale Image Quality Transformer~\cite{musiq} is a
\emph{no-reference} metric that predicts a human mean-opinion-score from the
reconstruction alone, without access to the original. Higher is better.
MUSIQ captures whether the codec's output looks plausible regardless of its
fidelity to the source.

\textbf{CLIP-IQA.} CLIP-IQA~\cite{clipiqa} uses a CLIP image encoder to
score reconstructions against natural-language quality prompts (``a good
photo'' versus ``a bad photo''). It is no-reference and semantically
grounded. Higher is better.

\textbf{FID.} Fr\'echet Inception Distance~\cite{fid} measures the
2-Wasserstein distance between the Inception-V3 feature statistics of the
reconstruction \emph{set} and the source \emph{set}. Lower is better. FID is
distributional rather than per-image: it captures whether the codec's
reconstructions, taken as a population, resemble natural photographs. The
original OSCAR paper reports FID only on CLIC2020 (its largest set), noting
that smaller sets do not give statistically stable FID estimates. Our
5000-image subset is more than ten times larger than CLIC2020 and provides a
stable FID measurement at every bit-rate.

\subsection{Evaluation Protocol}

For each of the eight bit-rate operating points $r \in \{0.0019,$ $0.0098,$
$0.0313,$ $0.0430,$ $0.0625,$ $0.0781,$ $0.0937,$ $0.125\}$, we feed every
image through the OSCAR forward pass of Section~\ref{sec:algorithm}, save the
reconstruction as an 8-bit PNG, and compute the six per-image metrics
(PSNR, MS-SSIM, DISTS, LPIPS, MUSIQ, CLIP-IQA) against the centre-cropped
source. The averages are taken with a NaN-skipping mean (Section~\ref{sec:implementation}).
FID is computed once per bit-rate by feeding the bit-rate's full 5000-image
reconstruction folder and the centre-cropped source folder to \texttt{pyiqa}'s
FID estimator, which extracts Inception-V3 features over both populations
and reports the resulting 2-Wasserstein distance.

% ----------------------------------------------------------------------------
\section{Problems Faced and Resolutions}
\label{sec:problems}

A faithful reproduction of OSCAR on AMD hardware ran into eight non-trivial
issues, ranging from a public mirror that disappeared mid-project to a
Windows power policy that silently froze an overnight run. We document each
below; many of the resolutions are reusable on related projects.

\subsection{Model Weight Acquisition}

The most immediate blocker was the disappearance of OSCAR's backbone. The
canonical \texttt{stabilityai/stable-diffusion-2-1} HuggingFace repository,
referenced verbatim by the upstream OSCAR README, no longer exists. Both
\texttt{snapshot\_download} and the diffusers single-file loader return
HTTP~404; an authenticated listing of the \texttt{stabilityai} organization
shows that none of the SD-2.x family remains hosted. The
\texttt{stabilityai/stable-diffusion-2} repository, used internally by
diffusers' \texttt{convert\_from\_ckpt.py} to fetch the CLIP-ViT-H text
encoder configuration, is also gone.

We reconstruct SD-2.1 from \texttt{webui/stable-diffusion-2-1}, a community
mirror that retains the original FP32 \texttt{safetensors} file. The
diffusers conversion utility is then redirected with three string
replacements in \texttt{convert\_from\_ckpt.py} from
\texttt{stabilityai/stable-diffusion-2} to
\texttt{sd-research/stable-diffusion-2-1-base}, an ungated mirror with a
byte-identical CLIP-ViT-H text encoder configuration. After conversion the
resulting diffusers folder is functionally equivalent to the canonical
release, and OSCAR's \texttt{AutoencoderKL.from\_pretrained} and
\texttt{UNet2DConditionModel.from\_pretrained} calls succeed unchanged.

\subsection{Porting to Windows ROCm}

\textbf{Distributed-training symbols missing.} PyTorch
\texttt{2.9.1+rocm7.2.1} on Windows is compiled with
\texttt{USE\_DISTRIBUTED=0}, so \texttt{torch.distributed.is\_available()}
returns \texttt{False} and most of the module surface is absent:
\texttt{is\_initialized}, \texttt{get\_world\_size}, \texttt{group},
\texttt{ReduceOp}, and the collective operations. Several libraries that
the OSCAR pipeline imports indirectly (most prominently
\texttt{vector-quantize-pytorch}) reference these symbols unconditionally at
module top level, causing an \texttt{ImportError} before any user code runs.

Our first attempt at a fix, \texttt{sitecustomize.py} placed in
\texttt{site-packages/} so that it would patch \texttt{torch.distributed} at
interpreter startup, caused a process fork bomb: every Python subprocess
spawned by the ROCm runtime itself re-triggered the sitecustomize file,
re-imported torch, and spawned more subprocesses recursively, reaching over
one hundred concurrent processes before we killed them. The final fix is a
manual-import compatibility module,
\texttt{diffusion/\_compat\_torch\_distributed.py}, that adds stubs
(\texttt{is\_initialized()=False}, \texttt{get\_world\_size()=1},
\texttt{group.WORLD=None}, and a \texttt{ReduceOp} class with the standard
enum members) when \texttt{is\_available()} is \texttt{False}. The module
is explicitly imported at the top of \texttt{main\_test.py} before any
library that touches \texttt{torch.distributed} is loaded.

\textbf{MSVC \texttt{MAX\_PATH} for \texttt{compressai}.}
\texttt{compressai 1.2.6} ships a \texttt{pybind11}-based ANS entropy coder
that must be built from source on Windows. The build fails with
\begin{quote}
\small\texttt{fatal error C1083: Cannot open compiler generated file: '': Invalid argument},
\end{quote}
the classic symptom of Windows' 260-character \texttt{MAX\_PATH} limit being
exceeded inside the \texttt{pip} build directory under
\texttt{C:{\textbackslash}Users{\textbackslash}User{\textbackslash}AppData{\textbackslash}Local{\textbackslash}Temp{\textbackslash}\ldots}.
The resolution is a single environment-variable override before the
\texttt{pip} invocation: setting \texttt{TMP=C:{\textbackslash}T} and
\texttt{TEMP=C:{\textbackslash}T} moves the build directory to a short
path and lets MSVC complete.

\textbf{\texttt{openai-clip}'s \texttt{pkg\_resources} import.}
\texttt{pyiqa}'s CLIP-IQA metric transitively imports
\texttt{openai-clip 1.0.1}, which begins with
\texttt{from pkg\_resources import packaging}. In modern
\texttt{setuptools} versions \texttt{packaging} is no longer re-exported
through \texttt{pkg\_resources}, causing an \texttt{ImportError}. A
three-line \texttt{try/except} around the import in
\texttt{site-packages/clip/clip.py} that falls back to
\texttt{import packaging} restores compatibility without otherwise touching
the third-party library.

\subsection{Numerical and Pipeline Robustness}

\textbf{DISTS produces NaN on CUDA FP32.} Running DISTS on the AMD GPU in
FP32 returns \texttt{NaN} for all inputs, including identical pairs (where
the correct answer is exactly zero). Probing the metric's internals
reveals that the second VGG stage of its feature extractor emits
\texttt{NaN} across the entire activation tensor, propagating downstream.
The same metric works correctly on the CPU and on the GPU in FP64; the
failure is specific to some FP32 operator kernel in this ROCm build. The
workaround is to instantiate DISTS with \texttt{device=``cpu''} and to
\texttt{.cpu()} the tensors at the call site. At roughly 80\,ms per pair
on the 9800X3D, the per-image overhead is acceptable.

\textbf{MS-SSIM crashes on small inputs.} An early overnight run halted
with
\begin{quote}
\small\texttt{RuntimeError: Calculated padded input size per channel: (8 x 16). Kernel size: (11 x 11). Kernel size can't be greater than actual input size}
\end{quote}
emitted from \texttt{pyiqa}'s MS-SSIM at image 345 of the lowest bit-rate
folder. The metric downsamples its input five times and then convolves an
$11 \times 11$ window over the result; the deepest scale needs each
spatial side to be at least 352~pixels for the kernel to fit. Roughly
70\,\% of the ImageNet validation images fall below this after the
upstream's divisible-by-128 trim. Two complementary fixes apply: the input
pipeline is standardized to $384 \times 384$
(Section~\ref{sec:implementation}), which satisfies the MS-SSIM constraint
with margin; and every \texttt{pyiqa} metric call is wrapped in a
\texttt{try/except} that returns \texttt{NaN} on failure, with per-rate
averaging via a NaN-skipping mean. Either fix alone would prevent
recurrence; both together make the run robust to any future
metric-specific crash on a single image.

\textbf{Mixed-precision was a no-op.} The upstream \texttt{main\_test.py}
defines a \texttt{weight\_dtype} variable but never applies it to the
model. We wrap the forward pass in
\texttt{torch.autocast(device\_type=``cuda'', dtype=weight\_dtype)}
(Section~\ref{sec:implementation}); the resulting 1.7$\times$ speedup is
what makes a full 5000-image run feasible in a single overnight slot.

\subsection{Operational Pitfalls}

\textbf{Auto-sleep silently stalled the eval.} The first full overnight
run produced only 345 reconstructions at the lowest bit-rate after seven
and a half elapsed hours, then resumed at nominal speed when we returned
to the machine. Inspection of the Windows event log identified a
\texttt{Kernel-Power Id=42, ``System Idle''} sleep event 14 minutes after
launch. The detached Python process was suspended along with the rest of
the system and resumed only when the display was re-engaged. The fix is
to disable system sleep before detached launches with
\texttt{powercfg /change standby-timeout-ac 0}. Once applied, a subsequent
5000-image run over all eight bit-rates completed in 4~hours~27~minutes
without any stalls.

% ----------------------------------------------------------------------------
\section{Results and Discussion}
\label{sec:results}

The full evaluation, comprising 5000 images at each of eight bit-rates,
completed in 4~hours~27~minutes on the 9070 XT (a steady 2.55 images per
second at FP16), with no metric failures and no stalls.
Table~\ref{tab:results} reports the per-bit-rate averages of all seven
metrics across the 5000 images.

\begin{table*}[t]
\caption{Reproduction results on the 5000-image ImageNet validation subset. Each row is a bit-rate operating point; columns are the seven metrics defined in Section~\ref{sec:methodology}. Up-arrows ($\uparrow$) indicate that higher is better; down-arrows ($\downarrow$) indicate that lower is better.}
\label{tab:results}
\begin{center}
\begin{tabular}{rrrrrrrr}
\toprule
$r$ (bpp) & PSNR $\uparrow$ & MS-SSIM $\uparrow$ & DISTS $\downarrow$ & LPIPS $\downarrow$ & MUSIQ $\uparrow$ & CLIP-IQA $\uparrow$ & FID $\downarrow$ \\
\midrule
0.0019 & 15.45 & 0.406 & 0.285 & 0.613 & 65.35 & 0.604 & 110.64 \\
0.0098 & 17.36 & 0.557 & 0.229 & 0.528 & 65.28 & 0.653 &  53.19 \\
0.0313 & 19.52 & 0.715 & 0.172 & 0.395 & 66.40 & 0.686 &  19.06 \\
0.0430 & 19.93 & 0.740 & 0.158 & 0.367 & 66.59 & 0.692 &  14.64 \\
0.0625 & 20.29 & 0.800 & 0.133 & 0.306 & 68.08 & 0.709 &   8.72 \\
0.0781 & 21.36 & 0.835 & 0.116 & 0.272 & 67.34 & 0.706 &   7.27 \\
0.0937 & 22.12 & 0.853 & 0.100 & 0.244 & 67.67 & 0.712 &   6.34 \\
0.1250 & 22.82 & 0.874 & 0.088 & 0.218 & 67.68 & 0.716 &   5.78 \\
\bottomrule
\end{tabular}
\end{center}
\end{table*}

\subsection{Quantitative Behaviour}

Every full-reference metric is monotone in the expected direction across
all eight bit-rates: PSNR rises from 15.45\,dB at 0.0019\,bpp to 22.82\,dB
at 0.125\,bpp; MS-SSIM rises from 0.406 to 0.874; DISTS falls from 0.285
to 0.088; LPIPS falls from 0.613 to 0.218. The two no-reference metrics
(MUSIQ and CLIP-IQA) move much less, rising only from roughly 65 to 68
and from 0.60 to 0.72 respectively. This pattern confirms the central
design claim of OSCAR: the bit-rate-to-timestep table behaves as a
single quality knob, with higher rates producing reconstructions that
are closer to the source on every reference-based measure.

The headline result is the FID drop, from 110.64 at 0.0019\,bpp to 5.78
at 0.125\,bpp. This is the widest dynamic range of any of the metrics
in our table, and arguably the metric most informative for a generative
codec: PSNR rewards pixel-exact matching that diffusion reconstructions
deliberately discard, but FID asks whether the codec's output looks like
a natural photograph at the population level. The lowest-bit-rate FID of
roughly 110 indicates that 0.0019\,bpp pushes the codec beyond the
diffusion prior's effective regime, producing recognizable but visibly
artificial reconstructions; by 0.0625\,bpp FID is already below 9, in the
range usually associated with high-fidelity generation.

The near-flat behaviour of MUSIQ and CLIP-IQA deserves a separate
comment. Both are no-reference metrics: they assess only whether the
reconstruction looks like a plausible photograph, irrespective of how
close it is to the source. The fact that they hover in the 65-68 and
0.60-0.72 ranges across the entire bit-rate ladder is the desirable
property: the diffusion prior produces output that looks natural at every
operating point, even when the underlying content is poorly recovered at
the lowest bit-rate. This is the property that distinguishes generative
codecs from PSNR-optimal learned codecs, whose low-bit-rate outputs look
blurred or blocked rather than incorrect but plausible.

\subsection{Comparison with the Original Paper}

The original OSCAR paper reports FID only on CLIC2020, with a range of
approximately 26 at the lowest bit-rate to 6 at the highest~\cite{oscar}.
Our ImageNet range of 110 to 6 is wider at the low end and matches at the
high end. The discrepancy at the low end is consistent with the
difference in test domains: CLIC2020 is a curated $1024 \times 1024$ set
of high-resolution natural photographs, while ImageNet validation
contains web-quality and heavily-cropped subjects of widely varying
composition. At extreme low bit-rates the diffusion prior has less to
work with on a generic ImageNet image than on a clean CLIC photo, which
inflates the low-bit-rate FID without affecting the high-bit-rate regime
where the codec carries enough information for the prior to be
secondary.

The per-image reference metrics show the same shape, if not the same
absolute magnitude, as the paper's CLIC2020 numbers. The paper reports
PSNR roughly in the range 19\,dB to 25\,dB across bit-rates; our ImageNet
results lie between 15.5\,dB and 22.8\,dB, with the low end depressed and
the high end close. This is again the expected outcome for a generative
codec evaluated on a more diverse domain.

\subsection{Throughput}

The 9070 XT processes a single $384 \times 384$ image through the full
OSCAR pipeline (VAE encode, hyper-encoder, U-Net, scheduler step, VAE
decode, six per-image metric evaluations, and a DISTS evaluation on the
CPU) at FP16 in roughly 0.4 seconds. Peak VRAM is 5.1\,GB at the lowest
bit-rate (largest pseudo-timestep) and stays below this for all other
rates. The total wall-clock of 4 hours 27 minutes compares favourably
with the multi-day runtime that would have resulted from running the
upstream code unmodified at varying ImageNet resolutions.

% ----------------------------------------------------------------------------
\section{Possible Improvements}
\label{sec:improvements}

Several directions could substantially improve OSCAR, both in compression
quality and in deployment cost.

\textbf{Tiled inference for high-resolution images.} The released
\texttt{main\_test.py} declares \texttt{vae\_decoder\_tiled\_size},
\texttt{vae\_encoder\_tiled\_size}, \texttt{latent\_tiled\_size}, and
\texttt{latent\_tiled\_overlap} arguments that are parsed but never wired
into the model. Implementing tiled VAE encode/decode and overlap-blended
U-Net inference would let OSCAR process arbitrary resolutions without the
fixed center-crop we adopt in this work, and would be especially valuable
on hardware with limited VRAM where a multi-megapixel image cannot fit in
a single forward pass.

\textbf{Continuous bit-rate via a learned mapping.} The bit-rate-to-timestep
mapping is a hard-coded eight-element table (Section~\ref{sec:algorithm}).
Replacing it with a small learned function $t(r)$, with $r$ continuous,
would allow OSCAR to interpolate between operating points and to support
rates that were not in the original training schedule. Combined with a
learnable hyper-encoder family parameterized by $r$, this would convert
OSCAR from an eight-rate codec into a true continuous-rate codec.

\textbf{Distillation of the U-Net.} The reconstruction cost is dominated
by a single forward pass through the 865M-parameter SD-2.1 U-Net. A
distilled or smaller U-Net, fine-tuned to reproduce the OSCAR
LoRA-adapted behaviour at each pseudo-timestep, could deliver an
order-of-magnitude speedup with limited quality loss. Recent work on
one-step diffusion distillation~\cite{vpred} is directly applicable.

\textbf{Joint hyper-encoder training.} OSCAR's stage-1 training fits the
eight hyper-encoders independently, which is convenient but ignores the
correlation between adjacent rates: a hyper-encoder for 0.0625\,bpp
shares most of its representational role with one for 0.0781\,bpp. A
joint training stage with a shared backbone and rate-conditioned heads
could reduce total parameter count and improve transfer between
operating points.

\textbf{End-to-end rate coding.} OSCAR's eight rates are realized through
discrete VQ configurations $(\text{cfg}_{ss}, \text{cfg}_{cs})$; the
actual bit-rate is approximately $\text{cfg}_{ss}^2 \cdot \log_2 \text{cfg}_{cs} / (H \cdot W)$.
A learned entropy model over the VQ indices, similar to the hyperprior
in mainstream learned codecs, would tighten this accounting and would
allow the hyper-encoder spatial size and codebook size to be jointly
optimized against rate at training time.

% ----------------------------------------------------------------------------
\section{Conclusion}
\label{sec:conclusion}

We have reproduced OSCAR, a one-step diffusion-based image codec, on the
ImageNet-1k validation set using consumer AMD hardware via ROCm 7.2.1. On
a 5000-image seed-controlled subset (more than an order of magnitude
larger than the paper's largest test set), the released checkpoint
reproduces the original result's central finding: a single shared
diffusion U-Net, conditioned on a bit-rate-mapped pseudo-timestep,
delivers monotone rate-distortion behaviour on all seven evaluation
metrics, with FID dropping from 110.64 at 0.0019\,bpp to 5.78 at
0.125\,bpp. The non-trivial porting work documented in
Section~\ref{sec:problems}, together with the input-pipeline and
metric-robustness modifications of Section~\ref{sec:implementation},
make the reproduction repeatable on hardware and at a dataset scale that
neither the original paper nor its public implementation contemplated.

% ----------------------------------------------------------------------------
\begin{thebibliography}{00}
\bibitem{oscar} J.~Guo \emph{et al.}, ``OSCAR: One-Step Diffusion Codec Across Multiple Bit-rates,'' in \emph{Advances in Neural Information Processing Systems (NeurIPS)}, 2025. arXiv:2505.16091.

\bibitem{hific} F.~Mentzer, G.~Toderici, M.~Tschannen, and E.~Agustsson, ``High-fidelity generative image compression,'' in \emph{Advances in Neural Information Processing Systems (NeurIPS)}, 2020.

\bibitem{ldm} R.~Rombach, A.~Blattmann, D.~Lorenz, P.~Esser, and B.~Ommer, ``High-resolution image synthesis with latent diffusion models,'' in \emph{IEEE/CVF Conference on Computer Vision and Pattern Recognition (CVPR)}, 2022, pp.~10684-10695.

\bibitem{perco} M.~Careil, M.~Muckley, J.~Verbeek, and S.~Lathuili\`{e}re, ``Towards image compression with perfect realism at ultra-low bitrates,'' in \emph{International Conference on Learning Representations (ICLR)}, 2024.

\bibitem{codiff} J.~Guo \emph{et al.}, ``CODiff: Conditional diffusion-based image compression.'' % TODO: verify exact title/venue/year from github.com/jp-guo/CODiff.

\bibitem{imagenet} O.~Russakovsky \emph{et al.}, ``ImageNet large scale visual recognition challenge,'' \emph{International Journal of Computer Vision (IJCV)}, vol.~115, no.~3, pp.~211-252, 2015.

\bibitem{ddpm} J.~Ho, A.~Jain, and P.~Abbeel, ``Denoising diffusion probabilistic models,'' in \emph{Advances in Neural Information Processing Systems (NeurIPS)}, 2020.

\bibitem{vpred} T.~Salimans and J.~Ho, ``Progressive distillation for fast sampling of diffusion models,'' in \emph{International Conference on Learning Representations (ICLR)}, 2022.

\bibitem{lora} E.~J.~Hu, Y.~Shen, P.~Wallis, Z.~Allen-Zhu, Y.~Li, S.~Wang, L.~Wang, and W.~Chen, ``LoRA: Low-rank adaptation of large language models,'' in \emph{International Conference on Learning Representations (ICLR)}, 2022.

\bibitem{msssim} Z.~Wang, E.~P.~Simoncelli, and A.~C.~Bovik, ``Multiscale structural similarity for image quality assessment,'' in \emph{Asilomar Conference on Signals, Systems, and Computers}, 2003, pp.~1398-1402.

\bibitem{lpips} R.~Zhang, P.~Isola, A.~A.~Efros, E.~Shechtman, and O.~Wang, ``The unreasonable effectiveness of deep features as a perceptual metric,'' in \emph{IEEE/CVF Conference on Computer Vision and Pattern Recognition (CVPR)}, 2018, pp.~586-595.

\bibitem{dists} K.~Ding, K.~Ma, S.~Wang, and E.~P.~Simoncelli, ``Image quality assessment: Unifying structure and texture similarity,'' \emph{IEEE Transactions on Pattern Analysis and Machine Intelligence (TPAMI)}, vol.~44, no.~5, pp.~2567-2581, 2022.

\bibitem{musiq} J.~Ke, Q.~Wang, Y.~Wang, P.~Milanfar, and F.~Yang, ``MUSIQ: Multi-scale image quality transformer,'' in \emph{IEEE/CVF International Conference on Computer Vision (ICCV)}, 2021, pp.~5148-5157.

\bibitem{clipiqa} J.~Wang, K.~C.~K.~Chan, and C.~C.~Loy, ``Exploring CLIP for assessing the look and feel of images,'' in \emph{AAAI Conference on Artificial Intelligence}, 2023.

\bibitem{fid} M.~Heusel, H.~Ramsauer, T.~Unterthiner, B.~Nessler, and S.~Hochreiter, ``GANs trained by a two time-scale update rule converge to a local Nash equilibrium,'' in \emph{Advances in Neural Information Processing Systems (NeurIPS)}, 2017.

\bibitem{pyiqa} C.~Chen and J.~Mo, ``IQA-PyTorch: PyTorch toolbox for image quality assessment,'' 2022. [Online]. Available: \url{https://github.com/chaofengc/IQA-PyTorch}

\end{thebibliography}

\end{document}
